Malé akademické oddelenia vykonávajúce výskum v oblasti strojového učenia čelia pretrvávajúcej infraštruktúrnej výzve. Nemôžu si dovoliť prístup s dvoma oddelenými systémami, bežný vo väčších inštitúciách, t. j. tradičný HPC plánovač ako Slurm pre dávkové tréningové úlohy popri samostatnom Kubernetes nasadení pre interaktívne služby, no prevádzka bez akejkoľvek orchestrácie vedie k neefektívnemu využívaniu zdrojov, manuálnym konfliktom v plánovaní a nízkej reprodukovateľnosti prostredí. Táto diplomová práca predstavuje návrh a implementáciu modulárneho Kubernetes-natívneho výpočtového klastra, ktorý konsoliduje dávkové úlohy, interaktívne vývojové prostredia a podporné služby na jedinom odľahčenom riadiacom uzle, navrhnutom pre akademické prostredia s obmedzenými zdrojmi.

Platforma je postavená na K3s, odľahčenej distribúcii Kubernetes, ktorej minimálna spotreba zdrojov umožňuje prevádzkovať produkčný riadiaci uzol na malom klastri spotrebiteľského hardvéru. Inteligentná vrstva plánovania založená na Volcano poskytuje spravodlivé zdieľanie zdrojov prostredníctvom front, gang scheduling pre distribuované úlohy a preempciu na základe priorít, čím kompenzuje absenciu hardvérového delenía GPU (NVIDIA MIG) na spotrebiteľských grafických kartách. Medzi doplnkové komponenty patrí JupyterHub pre interaktívny prístup k notebookom, Kube-Ray pre distribuovaný tréning, MLflow pre sledovanie experimentov a súkromný register kontajnerov pre reprodukovateľné prostredia.

Celý stack je navrhnutý pre prevádzku jediným správcom: bootstrap uzlov je automatizovaný prostredníctvom Ansible, životný cyklus aplikácií je spravovaný cez Helm chart-y nasadzované prostredníctvom CI/CD pipeline a celá konfigurácia je verzionovaná. Architektúra je evaluovaná na viac-uzlovom klastri heterogénneho spotrebiteľského hardvéru (NVIDIA Titan a GeForce RTX GPU, 12--24\,GB VRAM), čím sa demonštruje, že jediný zjednotujúci riadiaci uzol môže efektívne obslúžiť rôznorodé požiadavky na pracovné zaťaženie akademického laboratória pre výskum AI bez potreby dedikovaného DevOps tímu.
